\scanpage{16}\mspage{7}
d'où en ajoutant
\[
2\,x\wedge y=0\qquad\text{i.e.}\quad 2e_M=0
\]
ce qui impliquerait $2.1_\Lambda=0$, contrairement à l'hypothèse. On doit donc avoir
\[
\left\{\begin{aligned}
&x+y\deleted{\unsure{+}}\unsure{-|+}z=0\quad\text{ou}\quad x-y+z=0\\
&\text{i.e.}\quad z=x+y\quad\text{ou}\quad y=x+z
\end{aligned}\right.\tag*{(37)}
\]
Donc, l'ens des sommets de $SP_M$ incidents à $s\in S_M$, $s=\dot x$, s'identifie à la droite affine \inserted{$L_x$} sur $\Lambda$ formée des $y\in M$ satisfaisant (36), avec la structure de contour où les ``arêtes''
\[
\{y,z\}\in\unsure{\mathfrak P|\mathscr P}_2(L_x)\qquad L_x=\{y\in M\mid x\wedge y=e_M\}
\]
\deleted{sont \unsure{fo|for}} du contour sont formées des paires $\{y,z\}$ tels que l'on ait
\[
y-z\in\{\pm x\}.
\]
Les composantes connexes du contour correspondent \inserted{sont} alors \deleted{biunivoquement} \inserted{les \deleted{\unsure{ens|aux}} parties de} \deleted{aux} $L_x$ de la forme
\[
\{y+\unsure{n|k}x\mid \unsure{n|k}\in\mathbb{Z}\}=y+\Lambda_0 x\tag*{(38)}
\]
% archive sheet number: 15
